% Laget automatisk av elevark/lag_elevark.py fra sim/vogn-og-lodd.html. Endringer her blir overskrevet.
\documentclass[11pt]{article}
\usepackage{elevark}
\begin{document}
\elevarkhode{Fysikk 1 · Krefter}{Vogn trukket av lodd}{https://lektorbo.no/sim/vogn-og-lodd.html}{Simuleringen}

En vogn står på et bord og er festet i en snor som går over en trinse til et lodd. Når vi slipper, trekker loddet vogna med seg. Hvor stor blir akselerasjonen? Og er snorkraften like stor som tyngden til loddet?

\section*{Oppgaver}

Start med $M = 500\,\mathrm{g}$, $m = 100\,\mathrm{g}$ og ingen friksjon.

\begin{oppgaveliste}
\oppgave \fase{gjett}{Gjett}Er snorkraften større, mindre eller like stor som tyngden til loddet? Begrunn svaret før du ser på tallene.
\svar
\oppgave \fase{forklar}{Forklar}Tegn kreftene på vogna og på loddet hver for seg. Skriv opp Newtons 2. lov for hvert legeme, og vis at $a = mg / (M + m)$ uten friksjon.
\svar
\oppgave \fase{test}{Regn}Regn ut akselerasjonen og snorkraften. Hvor lang tid bruker vogna på 80~cm? Sjekk med simuleringen.
\svar
\oppgave \fase{test}{Undersøk}Gjør loddet tyngre og tyngre. Hva nærmer akselerasjonen seg? Hvorfor kan den aldri bli større enn $g$?
\svar
\oppgave \fase{test}{Regn}Sett friksjonen til $μ = 0{,}10$. Regn ut akselerasjonen og snorkraften. Hvor tungt må loddet være for at vogna i det hele tatt skal begynne å bevege seg?
\svar
\oppgave \fase{forklar}{Forklar}Hvorfor er det galt å si at «kraften som virker på vogna er tyngden til loddet»?
\svar
\oppgave \fase{utfordring}{Utfordring}I et forsøk måler dere akselerasjonen for mange ulike lodd, med samme totale masse $M + m$. Dere flytter lodd fra vogna til snora. Hva slags graf får dere av $a$ mot $m$? Hva er stigningstallet?
\svar
\end{oppgaveliste}

\end{document}
